\chapter{Tiền, nghề và tương lai còn mù mờ}
\markboth{Tiền, nghề và tương lai còn mù mờ}{Tiền, nghề và tương lai còn mù mờ}

\section{Nhà em cần tiền sớm, ước mơ để sau được không?}
\begin{truyen}{Nhà em cần tiền sớm, ước mơ để sau được không?}{Classroom Talk}
\chuhoa{H}{ọc sinh: ``Nhà em không dư giả. Em nghĩ kiếm tiền sớm mới là chuyện chính. Ước mơ để sau được không thầy?''}
\textit{Student: ``My family is not financially comfortable. I feel that earning money early is the main thing. Can dreams wait until later?''}\\[4pt]

\teacherpair{Được. Không phải ai cũng có đặc quyền mơ đẹp trước. Nhưng kể cả khi cần tiền sớm, em vẫn phải chọn cách kiếm tiền khiến mình lớn lên chứ không mòn đi.}{Yes, they can. Not everyone has the privilege of chasing beautiful dreams first. But even if you need money early, you still have to choose a way of earning that helps you grow instead of wearing you down.}

\studentpair{Nghĩa là sao?}{What do you mean by that?}

\teacherpair{Làm việc để sống là thật. Nhưng nếu chỉ nhìn cái trước mắt, em dễ bán rẻ những năm sau của mình.}{Working to survive is real. But if you only focus on what is in front of you right now, you may end up selling your later years too cheaply.}
\end{truyen}

\begin{baihoc}
Thực tế không giết ước mơ. Thực tế chỉ buộc mình chọn đường đi bền và tỉnh hơn.
\textit{(Reality does not kill dreams. It forces wiser and steadier paths toward them.)}
\end{baihoc}

\begin{ghinhoanh}
\textbf{Short English to remember:}

Urgency is real.

Do not let urgency steal your future.
\end{ghinhoanh}

\begin{tuhocnhanh}
1. Em đang cần tiền gấp đến mức nào? \textit{(How urgent is your financial reality?)}

2. Con đường nào vừa giúp em sống vừa giúp em tiến? \textit{(What path helps you survive and still move forward?)}
\end{tuhocnhanh}
\ngancach

\section{AI làm hết rồi, học sâu làm gì nữa}
\begin{truyen}{AI làm hết rồi, học sâu làm gì nữa}{Classroom Talk}
\chuhoa{H}{ọc sinh: ``Mai mốt AI làm hết, tụi em học cho lắm rồi cũng bị thay thôi. Vậy học sâu làm gì?''}
\textit{Student: ``If AI is going to do everything later anyway, then even if we study hard we may still be replaced. So why study deeply at all?''}\\[4pt]

\teacherpair{Nếu máy làm phần dễ hơn cho em, vậy em muốn thành người chỉ bấm nút hay thành người biết đặt câu hỏi, kiểm tra và quyết định?}{If machines handle the easier parts for you, then do you want to become a person who only presses buttons, or a person who knows how to ask questions, verify results, and make decisions?}

\studentpair{Nghe vậy thì học vẫn chưa được nghỉ ha thầy.}{Hearing it like that, I guess studying still does not get a break, right?}

\teacherpair{Đúng. Càng nhiều công cụ mạnh, người dùng càng phải bớt nông.}{Exactly. The more powerful the tools become, the less shallow the user can afford to be.}
\end{truyen}

\begin{baihoc}
Công cụ thông minh không cứu nổi người lười nghĩ.
\textit{(Smart tools cannot rescue lazy thinking.)}
\end{baihoc}

\begin{ghinhoanh}
\textbf{Short English to remember:}

Tools grow stronger.

Thinking must grow deeper.
\end{ghinhoanh}

\begin{tuhocnhanh}
1. Em muốn trở thành người dùng công cụ hay phụ thuộc công cụ? \textit{(Do you want to use tools, or depend on them?)}

2. Môn học nào đang rèn cho em cách nghĩ chứ không chỉ cách trả lời? \textit{(Which subject is training your thinking, not just your answers?)}
\end{tuhocnhanh}
\ngancach

\section{Em sợ chọn sai nghề rồi mắc kẹt cả đời}
\begin{truyen}{Em sợ chọn sai nghề rồi mắc kẹt cả đời}{Classroom Talk}
\chuhoa{H}{ọc sinh: ``Nghe người lớn hỏi chọn ngành là em mệt. Em sợ chọn sai một phát rồi kẹt luôn cả đời.''}
\textit{Student: ``I get tired every time adults ask me to choose a major. I am afraid of making one wrong choice and getting stuck for the rest of my life.''}\\[4pt]

\teacherpair{Sợ vậy là bình thường. Nhưng đời không chỉ có một nút bấm. Chọn nghề quan trọng, nhưng khả năng học lại và chỉnh hướng còn quan trọng hơn.}{That fear is normal. But life is not a machine with only one button. Choosing a career matters, but the ability to learn again and redirect yourself matters even more.}

\studentpair{Vậy sai vẫn cứu được hả thầy?}{So even a wrong choice can still be recovered from?}

\teacherpair{Sai có giá của nó. Nhưng nhiều người không chết vì chọn sai. Họ mòn vì chọn mà không chịu học tiếp.}{A wrong choice has its cost. But many people are not destroyed because they chose wrongly. They wear down because they choose something and then refuse to keep learning.}
\end{truyen}

\begin{baihoc}
Không ai đảm bảo chọn đúng tuyệt đối. Điều làm mình đỡ kẹt là năng lực đổi hướng khi cần.
\textit{(No one can guarantee a perfect first choice. The real protection is the ability to redirect.)}
\end{baihoc}

\begin{ghinhoanh}
\textbf{Short English to remember:}

No first choice is perfect.

Adaptability is part of career strength.
\end{ghinhoanh}

\begin{tuhocnhanh}
1. Em đang sợ chọn sai hay sợ chịu trách nhiệm sau khi chọn? \textit{(Do you fear the wrong choice, or the responsibility after choosing?)}

2. Kỹ năng nào giúp em đổi hướng nếu cần? \textit{(What skills would help you change direction later?)}
\end{tuhocnhanh}
\ngancach

\section{Không có quen biết thì cố có ăn thua không?}
\begin{truyen}{Không có quen biết thì cố có ăn thua không?}{Classroom Talk}
\chuhoa{H}{ọc sinh: ``Người ta có quan hệ, có chỗ đỡ đầu. Còn em không có ai. Vậy cố cho lắm có ăn thua không?''}
\textit{Student: ``Other people have connections and people backing them up. I have nobody. So does trying hard even make any difference?''}\\[4pt]

\teacherpair{Quan hệ là thật. Nhưng không có năng lực mà chỉ có quan hệ thì đi không xa. Còn có năng lực mà không biết mở miệng, không biết làm việc, cũng tự chặn đường mình.}{Connections are real. But if you have only connections and no ability, you will not go far. And if you have ability but do not know how to communicate or work with people, you block your own path too.}

\studentpair{Vậy em phải làm gì?}{So what should I do then?}

\teacherpair{Vừa học cho chắc, vừa tập cách cư xử, làm việc, giao tiếp cho đàng hoàng. Đó cũng là một dạng vốn.}{Study seriously, and at the same time learn how to behave, work, and communicate properly. That is also a form of capital.}
\end{truyen}

\begin{baihoc}
Đừng ngây thơ nghĩ chỉ cần giỏi là đủ. Cũng đừng cay cú tới mức quên xây giá trị của mình.
\textit{(Do not naively think skill alone solves everything. But do not become so bitter that you stop building your own value.)}
\end{baihoc}

\begin{ghinhoanh}
\textbf{Short English to remember:}

Connections matter.

So do trust and competence.
\end{ghinhoanh}

\begin{tuhocnhanh}
1. Em đang thiếu năng lực, thiếu cơ hội, hay thiếu cách kết nối? \textit{(Are you lacking skill, opportunity, or the ability to connect?)}

2. Em có đang bỏ quên kỹ năng làm việc với người khác không? \textit{(Are you neglecting the skill of working with others?)}
\end{tuhocnhanh}
\ngancach